\documentclass[11pt]{article}
\usepackage[margin=1in]{geometry}
\usepackage{amsmath,amssymb}
\usepackage[T1]{fontenc}
\usepackage{lmodern}
\usepackage{hyperref}

\title{Uniform Homeomorphism of Separable Banach Spaces: Known \texorpdfstring{$E_{\ell_\infty}$}{E ell-infinity} Lower Bound}
\author{Literature-implied partial status note}
\date{2026-06-16}

\begin{document}
\maketitle

\paragraph{Status.}
This is a \emph{literature-implied answer (partial subcase)} for the open problem in Ferenczi--Louveau--Rosendal, arXiv:math/0610289.  It is not a full solution of the analytic-completeness question.

\paragraph{Source question.}
In the section ``Uniform homeomorphism of complete separable metric spaces,'' after proving that uniform homeomorphism of complete separable metric spaces is complete analytic, Ferenczi--Louveau--Rosendal write that they have not been able to replace their metric spaces \(P(X)\) by Banach spaces and ask:
\[
\text{What is the complexity, with respect to }\leq_B,\text{ of uniform homeomorphism between separable Banach spaces?}
\]
In particular, they ask whether this relation is analytic complete.

\paragraph{Supporting theorem.}
Gao--Jackson--Sari, arXiv:0901.4092, explicitly identify this as the main problem left open from Ferenczi--Louveau--Rosendal.  Their Theorem 1.1 constructs a Borel family \(\{S_{\vec x}:\vec x\in\mathbb R^\omega\}\) of separable Banach spaces such that, for all \(\vec x,\vec y\),
\[
\vec x-\vec y\in \ell_\infty
\quad\Longleftrightarrow\quad
S_{\vec x}\text{ and }S_{\vec y}\text{ are uniformly homeomorphic}
\quad\Longleftrightarrow\quad
S_{\vec x}\text{ and }S_{\vec y}\text{ are locally equivalent}.
\]
Equivalently,
\[
E_{\ell_\infty}\leq_B E_{\mathrm{UH}}^{\mathrm{Banach}},
\]
where \(E_{\mathrm{UH}}^{\mathrm{Banach}}\) denotes uniform homeomorphism between separable Banach spaces.  Their Theorem 1.2 further proves that local equivalence between separable Banach spaces is Borel bireducible with \(E_{\ell_\infty}\).

\paragraph{Interpretation.}
This gives a nontrivial partial answer to the source problem: uniform homeomorphism of separable Banach spaces is at least complete \(K_\sigma\) in the Borel-reducibility hierarchy.  The same paper explains that this lower bound is the strongest result obtainable from local finite-dimensional structure alone, via Ribe's theorem and local equivalence.

\paragraph{Remaining open scope.}
The analytic-completeness part of the Ferenczi--Louveau--Rosendal problem is not answered by Gao--Jackson--Sari.  They state that proving \(E_0^\omega\leq_B E_{\mathrm{UH}}^{\mathrm{Banach}}\) would already show that uniform homeomorphism is strictly more complex than \(E_{\ell_\infty}\) and would move beyond the local-structure barrier.

\begin{thebibliography}{9}
\bibitem{FLR}
V. Ferenczi, A. Louveau, and C. Rosendal,
\emph{The complexity of classifying separable Banach spaces up to isomorphism},
arXiv:math/0610289.

\bibitem{GJS}
S. Gao, S. Jackson, and B. Sari,
\emph{On the complexity of the uniform homeomorphism relation between separable Banach spaces},
arXiv:0901.4092.
\end{thebibliography}

\end{document}

